\scanpage{80}
\[
 B_\Delta\simeq A_\Delta[T]/(T^2+bT+c)\qquad b,c\in A_\Delta
\]
avec l'idéal discriminant de $B_\Delta$ égal :
\[
 \unsure{\mathfrak{D}|\mathscr D}(B_\Delta/A_\Delta)=(b^2-4c)A_\Delta
\]
donc, à cause de l'hypothèse $B_\Delta$ étale sur $A_\Delta$
\[
 b^2-4c\in A_\Delta^*
\]
D'ailleurs, prenant la restriction à $A_{2\Delta}$, on trouve
\[
 b^2-4c=u^2\Delta,\qquad u\in A_{2\Delta}^*=(A_\Delta)_2^*
\]

\deleted{Supposons que}
\[
 \deleted{\boxed{(A_{2\Delta})_{\unsure{2|*}}^*=\{2^\mu\Delta^\nu\unsure{v|\upsilon}\mid\mu,\nu\in\mathbb Z,\ v\in A^*)}}\ \deleted{\text{\textcircled{\unsure{e}}}}
\]
\deleted{donc on aura}
\[
 \deleted{b^2-4c=2^{2\mu}\Delta^{2\nu+1}v^2\quad\text{dans }(A_{\deleted{\unsure{2}}\Delta})_2}
\]
\deleted{avec}
\[
 \deleted{v\in A^*}
\]

Comme $b^2-4c\in A_\Delta^*$, on conclut \deleted{que}\inserted{$u^2\in A_\Delta^*$} ce qui implique
\[
 \deleted{\text{Si }\boxed{2\ \text{non inv. dans } A_\Delta}\ \text{\textcircled{f}}\ \unsure{\Longleftarrow}}
\]
$*\ u\in A_\Delta^*$ si $\boxed{A_\Delta\ \text{intégralemt clos dans } A_{2\Delta}}$ \textcircled{f}

$\deleted{\mu=0}$, donc,
\[
 b^2-4c=\Delta^{\deleted{\unsure{2\nu+1}}}u^2\qquad u\in A^* \tag*{(\unsure{*|$\alpha$})}
\]
Inversement, si on a une relation \unsure{$\alpha$|$(*)$} (en \deleted{y supposant seulement $v\in A_\Delta^*$, et oubliant les hypothèses préliminaires faites}) l'on \deleted{en aura}
\[
 \deleted{b^2-4c=\lambda^2\Delta\qquad\lambda\in A^*\quad(\lambda=(\Delta^\nu v)^{2})}
\]

% archive sheet number: 79
